%!TEX encoding=UTF-8 Unicode
%!TEX encoding=UTF-8 Unicode
%tweeks on pdf version so everybody is happy
\directlua{pdf.setminorversion(4)} % for facile.cines.fr
%\pdfcompresslevel=0 % Not needed

\documentclass[xcolor={usenames,dvipsnames},hyperref={pdfusetitle}]{beamer}
\usepackage{lmodern}

%=========================Language and encoding ==============================

\usepackage[utf8]{inputenc}
\usepackage[english]{babel}
\usepackage[T1]{fontenc}
% Fix size errors due to T1 in bbl file
\usepackage{fix-cm}
%=============================================================================

%========================= Todo notes  =======================================

%!TEX encoding=UTF-8 Unicode

% Todo notes
\newcommand{\DB}[1]{\todo[author=David,inline]{#1}}
\newcommand{\GH}[1]{\todo[author=Guillaume,inline]{#1}}

% Usefull stuff
\newcommand{\Input}[1]{\input{tex/#1}}

% References
\newcommand{\fig}[1]{Figure~\ref{fig:#1}}
\newcommand{\tbl}[1]{Table~\ref{tab:#1}}
\newcommand{\alg}[1]{Algorithm~\ref{alg:#1}}
\newcommand{\sect}[1]{Section~\ref{sec:#1}}
\newcommand{\chap}[1]{Chapter~\ref{chap:#1}}

% Tools names
\newcommand{\Tabarnac}{\emph{Tabarnac}\xspace}
\newcommand{\Moca}{\emph{Moca}\xspace}


\usepackage{todonotes}
\presetkeys{todonotes}{inline}{}

%=============================================================================

%========================= Figures ===========================================

\usepackage[]{caption}
\usepackage[]{subcaption}
\usepackage{graphicx} % support the \includegraphics command and options
\graphicspath{{./img/}{./style/}{./tikz/}}
%!TEX encoding=UTF-8 Unicode

% Libraries
\usetikzlibrary{shapes,arrows,decorations,decorations.pathreplacing,decorations.markings,fit}
\usetikzlibrary{positioning,backgrounds,calc,patterns}



%!TEX encoding=UTF-8 Unicode
\usepackage{algorithm}
\usepackage{algpseudocode}
\usepackage{varwidth} % for algorithm in tikz node

%\algblockdefx[]{Function}{EndFunction}
%[2]{\algorithmicfunction\ \textproc{#1}{(#2)}}%
%{\algorithmicend\ \algorithmicfunction}%
%\algnewcommand\Callp[2]{\textproc{#1}(#2)}%

\algblockdefx[Program]{Program}{EndProgram}%
    [1]{\texttt{#1:}}%
    {}

\usepackage{listings}

%\renewcommand*\thelstnumber{\arabic{lstnumber}}

\DeclareCaptionFormat{mylst}{\hrule\vspace{1pt}#1#2#3}
\captionsetup[lstlisting]{format=mylst,labelfont=bf,singlelinecheck=off,labelsep=space}
%
\lstset{
    language=C,
    basicstyle=\small\ttfamily,
    numbers=left,
    numbersep=10pt,
    xleftmargin=20pt,
    frame=tb,
    framexleftmargin=20pt,
    escapeinside={|}{|},          % if you want to add LaTeX within your code
    keywordstyle=\bf,
}

% \lstset{ %
%   %backgroundcolor=\color{white},   % choose the background color; you must add \usepackage{color} or \usepackage{xcolor}
%   %basicstyle=\footnotesize,        % the size of the fonts that are used for the code
%   breakatwhitespace=false,         % sets if automatic breaks should only happen at whitespace
%   breaklines=true,                 % sets automatic line breaking
%   captionpos=t,                    % sets the caption-position to bottom
%   %commentstyle=\color{mygreen},    % comment style
%   %deletekeywords={...},            % if you want to delete keywords from the given language
%   escapeinside={|}{|},          % if you want to add LaTeX within your code
%   extendedchars=true,              % lets you use non-ASCII characters; for 8-bits encodings only, does not work with UTF-8
%   frame=tb,	                   % adds a frame around the code
%   keepspaces=true,                 % keeps spaces in text, useful for keeping indentation of code (possibly needs columns=flexible)
%   %keywordstyle=\color{blue},       % keyword style
%   language=C,                 % the language of the code
%   %otherkeywords={*,...},           % if you want to add more keywords to the set
%   numbers=left,                    % where to put the line-numbers; possible values are (none, left, right)
%   numbersep=5pt,                   % how far the line-numbers are from the code
%   numberstyle=\tiny, % the style that is used for the line-numbers
%   rulecolor=\color{black},         % if not set, the frame-color may be changed on line-breaks within not-black text (e.g. comments (green here))
%   showspaces=false,                % show spaces everywhere adding particular underscores; it overrides 'showstringspaces'
%   showstringspaces=false,          % underline spaces within strings only
%   showtabs=false,                  % show tabs within strings adding particular underscores
%   stepnumber=2,                    % the step between two line-numbers. If it's 1, each line will be numbered
%   stringstyle=\color{mymauve},     % string literal style
%   tabsize=2,	                   % sets default tabsize to 2 spaces
%   title=\lstname                   % show the filename of files included with \lstinputlisting; also try caption instead of title
% }

\usepackage{epstopdf}
\usepackage{booktabs}
\usepackage{multirow}
%\usepackage{subcaption}

%=============================================================================

%=============================================================================

%========================= Hyperref ==========================================


\hypersetup{
    colorlinks=true, %colore les liens
    breaklinks=true, %permet le retour à la ligne dans les liens trop longs
    urlcolor= blue, %couleur des hyperliens
    linkcolor= black, %couleur des liens internes
    bookmarksopen=true,
    citecolor=black,
}
%=============================================================================

%========================= Other useful includes =============================

\usepackage{ifthen}
\usepackage[absolute,overlay]{textpos} %to set some blocks position
%=============================================================================

%========================= Beamer theme =====================================

%Stuff for printable version
\mode<handout>{
    \usetheme{default}
    \setbeamercolor{background canvas}{bg=black!5}
    \pgfpagesuselayout{4 on 1}[a4paper,landscape,border shrink=2.5mm]
}

\usetheme{AntibesCompact}

\newcommand{\alertitem}{\item<+-|alert@+>}
\newcommand\alertblockat[3]{
    \alt#1{
        \begin{alertblock}{#2}
            #3
        \end{alertblock}
    }{
        \begin{block}{#2}
            #3
        \end{block}
    }
}

%=============================================================================

%========================= Title frame  ======================================
\title[]{Analyzing the memory behavior of parallel scientific applications}
\author[David Beniamine]{\textbf{David Beniamine}}
\institute[Polaris / Datamove]{
    \includegraphics[height=.10\textheight]{img/logoUGA.jpg}
    \quad
    \includegraphics[width=.10\textwidth]{img/LIG_coul.jpg}
    \quad
    \includegraphics[width=.15\textwidth]{img/inria.jpg}
    \quad
    \includegraphics[width=.18\textwidth]{img/polaris.png}
    \quad
    \includegraphics[width=.18\textwidth]{img/datamove.png}
}


\newcommand{\enableTocAtSection}{
    \AtBeginSection[]
    {
        \ifthenelse{\boolean{sectiontoc}}{
            \begin{frame}<beamer>
                \frametitle{Outline}
                \tableofcontents[currentsection,currentsubsection]
            \end{frame}
        }
    }
    \AtBeginSubsection[]
    {
        \ifthenelse{\boolean{sectiontoc}}{
            \begin{frame}<beamer>
                \frametitle{Outline}
                \tableofcontents[currentsection,currentsubsection]
            \end{frame}
        }
    }
}

%=============================================================================

\begin{document}
%========================= Title and outlines ================================
\begin{frame}{}
    \titlepage
\small
{\centering\itshape Jury members\par}
\begin{tabular}[t]{@{}l@{\hspace{3pt}}p{.45\textwidth}@{}}
President: & Pr, Martin Quinson\\
Reviewers: & Pr, Jes\'us Labarta Mancho \\
& Pr,  Raymond Namyst \\
Examiners: & Dr, Lucas M. Schnorr \\
\end{tabular}%
\begin{tabular}[t]{@{}l@{\hspace{3pt}}p{.45\textwidth}@{}}
Supervisors: & Pr, Bruno Raffin \\
 & Dr, Guillaume Huard
\end{tabular}%
\end{frame}

\newboolean{sectiontoc}
\setboolean{sectiontoc}{true} % default to true

%=============================================================================

%========================= Real presentation =================================

\section{Context}

\begin{frame}{Scientific applications and performances}
    \todo{Put something here}
\end{frame}

\begin{frame}{Improving sequential computer performances}
    \begin{columns}
        \begin{column}{.45\textwidth}
            \centering
            %!TEX encoding=UTF-8 Unicode
%Palette GnBu 5 col + white
\definecolor{ColPU}{HTML}{FFFFFF}
\definecolor{ColCore}{HTML}{F0F9E8}
\definecolor{ColL1}{HTML}{BAE4BC}
\definecolor{ColL2}{HTML}{7BCCC4}
\definecolor{ColL3}{HTML}{43A2CA}
\definecolor{ColM}{HTML}{0868AC}
\definecolor{ColS}{HTML}{FFFFFF}

\pgfdeclarelayer{bg}
\pgfdeclarelayer{bbg}
\pgfdeclarelayer{bbbg}
\pgfsetlayers{bbbg,bbg,bg,main}


\tikzset{
    box/.style={
        shape=rectangle,
        text centered,
        draw,
    },
}

\begin{tikzpicture}[font=\small, every pic/.style={scale=.9}]
    \node[box,fill=ColPU] (PU-0) at (0,0) {Thread};
    \node[minimum width=3.3em] (name) at ($(PU-0)+(0,1)$) {Core};

    \begin{pgfonlayer}{bg}
        \node[box,fill=ColCore, fit=(name) (PU-0) ] (core-0) {};
    \end{pgfonlayer}

    \node (cache) at ($(core-0.north)+(0,1)$) {};

    \uncover<4->{
        \draw (core-0.north) -- (cache);

        \draw[fill=ColL3] ($(core-0.north west)+(0,.5)$) rectangle
            ($(core-0.north east)+(0,1)$) node[pos=.5] {Cache};
    }

    \node (cpu-name) at ($(cache)+(0,.5)$) {CPU};

    \begin{pgfonlayer}{bbg}
        \node[box,fill=ColS,fit=(core-0) (cpu-name)] (cpu) {};
    \end{pgfonlayer}

    \draw[line width=.5em] (cpu.north) -- ($(cpu.north)+(0,1)$);

    \draw[fill=ColM, text=white] ($(cpu.north west)+(-.5,1)$) rectangle
        ($(cpu.north east)+(.5,2)$) node[pos=.5]{Memory};


\end{tikzpicture}
% vim: et si sta lbr  sw=4 ts=4 spelllang=en_us

        \end{column}
        \begin{column}{.45\textwidth}
            \pause
            \begin{block}{Increase cpu frequency}
                \begin{itemize}[<+->]
                    \item Done for years
                    \alertitem Gap CPU / Memory
                    \begin{itemize}
                        \alertitem Add memory caches
                    \end{itemize}
                    \alertitem Physical limits
                        \begin{itemize}
                            \alertitem Energy consumption
                            \alertitem Heat dissipation
                        \end{itemize}
                \end{itemize}
            \end{block}
            \uncover<8->{
                \begin{alertblock}{Build parallel machines}
                    Writing efficient code becomes complex
                \end{alertblock}
            }
        \end{column}
    \end{columns}
\end{frame}

\begin{frame}{An simple example}
    \centering
    \scalebox{.6}{
        %!TEX encoding=UTF-8 Unicode

\definecolor{ColI}{HTML}{1B9E77}
\definecolor{ColK}{HTML}{D95F02}
\definecolor{ColJ}{HTML}{7570B3}

\pgfdeclarelayer{background}
\pgfdeclarelayer{foreground}
\pgfsetlayers{background,foreground}


\tikzstyle{PrimaryA}   = [-latex,very thick]
\tikzstyle{SecondaryA} = [-latex,very thick,dashed]

\newcommand{\coli}[1]{\textcolor{ColI}{#1}}
\newcommand{\colj}[1]{\textcolor{ColJ}{#1}}
\newcommand{\colk}[1]{\textcolor{ColK}{#1}}

% #1: name, #2: start pos, #3: size
\newcommand{\matgrid}[3]{
        \draw #2 grid ($#2+(#3,#3)$);
        % Four corners
        \coordinate (m#1-00) at ($#2+(0.5,0.5)$);
        \coordinate (m#1-0N) at ($#2+(0.5,#3-0.5)$);
        \coordinate (m#1-N0) at ($#2+(#3-0.5,0.5)$);
        \coordinate (m#1-NN) at ($#2+(#3-0.5,#3-0.5)$);

        \node (#1) at ($(m#1-00)+(-1,0)$){#1};

        \node (#1) at ($(m#1-0N)+(-.2,.2)$)  {0};
        \node (#1) at ($(m#1-00)+(0,-.2)$) {N-1};
        \node (#1) at ($(m#1-NN)+(.2,0)$)  {N-1};
}


\begin{tikzpicture}[font=\small]

    \begin{pgfonlayer}{background}
        \node[draw,rounded corners] at (2.5,9){%
            \begin{varwidth}{\linewidth}
                \begin{algorithmic}
                    \For{\coli{i in 0..N-1}}
                    \For{\colj{j in 0..N-1}}
                    \For{\colk{k in 0..N-1}}
                    \State C[\coli{i},\colj{j}] += A[\coli{i},\colk{k}] * B[\colk{k},\colj{j}]
                            \EndFor
                        \EndFor
                    \EndFor
                \end{algorithmic}%
            \end{varwidth}%
        };

        \matgrid{A}{(0,0)}{5}
        \matgrid{B}{(6,6)}{5}
        \matgrid{C}{(6,0)}{5}

    \end{pgfonlayer}

    % Indexes

    \begin{pgfonlayer}{foreground}
        %% A
        \draw[PrimaryA,ColI]   (mA-0N) -- node [above] {i} (mA-NN);
        \draw[SecondaryA,ColK] (mA-0N) -- node [left]  {k} (mA-00);

        %% B
        \draw[PrimaryA,ColJ]   (mB-0N) -- node [left]  {j} (mB-00);
        \draw[SecondaryA,ColK] (mB-0N) -- node [above] {k} (mB-NN);

        %% C
        \draw[PrimaryA,ColI]   (mC-0N) -- node [above] {i} (mC-NN);
        \draw[SecondaryA,ColJ] (mC-0N) -- node [left]  {j} (mC-00);
    \end{pgfonlayer}

\end{tikzpicture}
% vim: et si sta lbr  sw=4 ts=4 spelllang=en_us

    }
\end{frame}

\begin{frame}{Parallel and NUMA machines}
    \centering
    \scalebox{.6}{
        %!TEX encoding=UTF-8 Unicode
%Palette GnBu 5 col + white
\definecolor{ColPU}{HTML}{FFFFFF}
\definecolor{ColCore}{HTML}{F0F9E8}
\definecolor{ColL1}{HTML}{BAE4BC}
\definecolor{ColL2}{HTML}{7BCCC4}
\definecolor{ColL3}{HTML}{43A2CA}
\definecolor{ColM}{HTML}{0868AC}
\definecolor{ColS}{HTML}{FFFFFF}

\pgfdeclarelayer{bg}
\pgfdeclarelayer{bbg}
\pgfdeclarelayer{bbbg}
\pgfsetlayers{bbbg,bbg,bg,main}


\tikzset{
    box/.style={
        shape=rectangle,
        text centered,
        draw,
    },
    pics/core/.style args={#1#2#3#4}{
        % Args: #1: nb PU, #2 core id, #3 direction: + or -
        code={
            \pgfmathtruncatemacro{\pmin}{#1*#2}
            \pgfmathtruncatemacro{\pmax}{\pmin+#1-1}
            %PUs
            \foreach \pu in {\pmin,...,\pmax}{
                \pgfmathtruncatemacro{\pstep}{\pu-\pmin}
                \node[box,fill=ColPU] (PU-\pu)at (0,#3.7*\pstep) {PU\#\pu};
            }
            %Core ID
            \node (inv) at ($(PU-\pmin)#3(0,-.3)$) {};
            \node[minimum width=3.3em] (name) at ($(PU-\pmin)!0.5!(PU-\pmax)#3(0,1)$) {Core\##2};
            % L1
            \begin{pgfonlayer}{bg}
                \node[box,fill=ColCore,inner sep=.1pt, fit=(name) (PU-\pmin)
                (PU-\pmax) (inv)] (core-#2) {};
            \end{pgfonlayer}
            \draw[fill=ColL1] ($(core-#2.#4 west)#3(0,.1)$) rectangle ($(core-#2.#4 east)#3(0,.5)$)%
            node[pos=.5] (l1) {L1};
            % links
            \draw (core-#2.#4) -- ($(core-#2.#4)#3(0,.1)$);
            \coordinate (l1-#2-n) at ($(core-#2.#4)#3(0,.5)$);
        }
    },
    pics/l2group/.style args={#1#2#3#4}{
        % Args: #1: nb Cores, #2 group id, #3 direction: + or -
        code={
            \pgfmathtruncatemacro{\cmin}{#1*#2}
            \pgfmathtruncatemacro{\cmax}{\cmin+#1-1}
            % Cores
            \foreach \core in {\cmin,...,\cmax}{
                \pgfmathtruncatemacro{\cstep}{\core-\cmin}
                \draw (1.4*\cstep,0) pic[font=\tiny] {core={2}{\core}{#3}{#4}};
            }
            % L2
            \draw[fill=ColL2] ($(core-\cmin.#4 west)#3(0,.6)$) rectangle
            ($(core-\cmax.#4 east)#3(0,1.1)$) node[pos=.5]{L2};
            % Coordinates for L3
            \coordinate (l2g-#2-w) at ($(core-\cmin.#4 west)#3(0,1.1)$);
            \coordinate (l2g-#2-e) at ($(core-\cmax.#4 east)#3(0,1.1)$);
            %links
            \foreach \core in {\cmin,...,\cmax}{
                \draw (l1-\core-n) -- ($(core-\core.#4)#3(0,.6)$);
            }
        }
    },
    pics/socket/.style args={#1#2#3#4}{
        % Args: #1: nb l2 groups, #2 socket id, #3 arguments for underlying pics
        code={
            \pgfmathtruncatemacro{\lmin}{#1*#2}
            \pgfmathtruncatemacro{\lmax}{\lmin+#1-1}
            % Cores
            \foreach \lg in {\lmin,...,\lmax}{
                \pgfmathtruncatemacro{\lgstep}{\lg-\lmin}
                \draw (2.8*\lgstep,0) pic[font=\small] {l2group={2}{\lg}{#3}{#4}};
            }
           % L3
            \draw[fill=ColL3] ($(l2g-\lmin-w)#3(0,.1)$) rectangle ($(l2g-\lmax-e)#3(0,.5)$)%
            node[pos=.5]{L3};
           % Coordinates for Mem
            \coordinate (s-#2-w) at ($(l2g-\lmin-w)#3(0,.5)$);
            \coordinate (s-#2-e) at ($(l2g-\lmax-e)#3(0,.5)$);
           % links
           %\foreach \lg in {\lmin,...,\lmax}{
           %    \draw ($(l2g-\lg-w)!.5!(l2g-\lg-e)$) --
           %        ($(l2g-\lg-w)!.5!(l2g-\lg-e)#3(0,.1)$);
           %}
           % CPU
            \ifthenelse{\equal{#3}{+}}{
                \node (minnode) at (-.6,-.3) {};
            }{
                \node (minnode) at (-.6,.3) {};
            }
            \node (maxnode) at ($(l2g-\lmax-e)#3(0,1)$) {};
            \node (sockname) at ($(l2g-\lmin-e)#3(0,.8)$) {Socket \##2};

            \begin{pgfonlayer}{bbg}
                \node[box,fill=ColS,fit=(minnode) (maxnode)] (cpu-#2) {};
            \end{pgfonlayer}
           % Memory
            \draw[fill=ColM,text=white] ($(cpu-#2.#4 west)#3(0,.5)$) rectangle ($(cpu-#2.#4 east)#3(0,1.5)$)%
            node[pos=.5]{Memory bank \##2};

            \draw[very thick] (cpu-#2.#4) -- ($(cpu-#2.#4)#3(0,.5)$);

        }
    },
}

\begin{tikzpicture}[font=\small, every pic/.style={scale=.9}]
    \pic at (0,0)  {socket={2}{0}{+}{north}};
    \pic at (6.5,0)  {socket={2}{1}{+}{north}};
    \pic at (0,-2) {socket={2}{2}{-}{south}};
    \pic at (6.5,-2) {socket={2}{3}{-}{south}};

    \begin{pgfonlayer}{bbbg}
        \draw[line width=1em] (cpu-0) -- (cpu-2);
        \draw[line width=1em] (cpu-0) -- (cpu-1);
        \draw[line width=.3em]  (cpu-0) -- (cpu-3);
        \draw[line width=1em] (cpu-1) -- (cpu-3);
        \draw[line width=1em] (cpu-2) -- (cpu-3);
        \draw[line width=.3em]  (cpu-1) -- (cpu-2);
    \end{pgfonlayer}

\end{tikzpicture}
% vim: et si sta lbr  sw=4 ts=4 spelllang=en_us

    }
    \pause
\end{frame}

\begin{frame}{Research statement}
    \begin{alertblock}{Statement}
        How can we analyze the memory behavior of an application to optimize its performances ?
    \end{alertblock}
    \pause
    \begin{block}{Challenges}
        \begin{itemize}[<+->]
                \alertitem Collect memory traces
                \begin{itemize}
                    \item \textbf{Complete}: do not miss part of the address space
                    \item \textbf{Precise}:  enough to detect patterns
                    \item \textbf{Detailed}: embed all meta data about the access
                \end{itemize}
                \alertitem Visualize memory traces
                \alertitem Take advantage of the obtained knowledge
        \end{itemize}
    \end{block}
\end{frame}

\enableTocAtSection
\begin{frame}{Outline}
    \tableofcontents
\end{frame}

\section{State of the art}

\begin{frame}{Generic performance analysis tools}
    <++>
\end{frame}

\begin{frame}{Existing memory profiling tools}
    <++>
\end{frame}

\section{Analyzing the global memory behavior}

\subsection{Tabarnac}

\begin{frame}{Trace collection}
    <++>
\end{frame}

\begin{frame}{Visualizations}
    <++>
\end{frame}

\subsection{Using Tabarnac to optimize a well known benchmark: IS}

\begin{frame}{First analysis}
    <++>
\end{frame}

\begin{frame}{Issue and optimization}
    <++>
\end{frame}

\begin{frame}{Second analysis}
    <++>
\end{frame}

\begin{frame}{Results}
    <++>
\end{frame}

\begin{frame}{Conclusions}
    <++>
\end{frame}

\section{Collecting and analyzing fine grained traces}

\subsection{Moca an efficient memory trace collection system}

\begin{frame}{Background knowledge}
    <++>
\end{frame}

\begin{frame}{Design}
    <++>
\end{frame}

\begin{frame}{Handling page faults}
    <++>
\end{frame}


\begin{frame}{Experimental setup}
    <++>
\end{frame}

\begin{frame}{Experimental results: trace precision}
    <++>
\end{frame}

\begin{frame}{Experimental results: overhead}
    <++>
\end{frame}

\subsection{Visualizing and analyzing Moca traces}

\begin{frame}{Framesoc / Ocelotl}
    <++>
\end{frame}

\begin{frame}{Example: Matrix multiplication}
    <++>
\end{frame}

\begin{frame}{Limits}
    <++>
\end{frame}

\begin{frame}{Programmatic approach}
    <++>
\end{frame}

\begin{frame}{Example: ??}
    <++>
\end{frame}

\section{Conclusions and perspectives}

\begin{frame}{Global overview}
    \cite{Beniamine15TABARNAC}
    <++>
\end{frame}

\begin{frame}{Fine trace collecion}
    \cite{Beniamine16Moca}
    <++>
\end{frame}

\begin{frame}{Trace analysis}
    <++>
\end{frame}

\begin{frame}{Perspectives}
    <++>
\end{frame}

\newcounter{finalframe}
\setcounter{finalframe}{\value{framenumber}}
%Last numbered frame go here

\begin{frame}{Perspectives}
\end{frame}
%=============================================================================

%=============================================================================
%Uncomment next lines for uncounted backup slides & biblio
\section*{Bibliography}
%
\bibliographystyle{apalike}
\bibliography{biblio}

%========================= Backup slides =====================================
%\section*{Hidden slides}
%put this line before each frame

%\setcounter{framenumber}{\value{finalframe}}
%=============================================================================
\end{document}




